% !PRINTED 34
\weberpage{34}
\begin{center}
2. \gesperrt{Die Landensche Transformation.}
\end{center}

\[
V = a\eta(\xi-\eta), \qquad U = \xi(\xi-\varkappa^2\eta).
\]

Die Bedingungen, daß $U-V$, $U-\lambda^2V$ Quadrate sind, lauten

\[
4a = (\varkappa^2+a)^2, \qquad 4a\lambda^2 = (\varkappa^2+a\lambda^2)^2,
\]

woraus

\[
\lambda(\varkappa^2+a) = \varkappa^2+a\lambda^2,
\]

oder durch $\lambda-1$ dividiert

\[
a\lambda = \varkappa^2;
\]

dies in eine der obigen Gleichungen eingesetzt, gibt

\[
4\lambda = \varkappa^2(1+\lambda)^2,
\]

und durch Auflösung dieser quadratischen Gleichung, wenn
$\varkappa'=\sqrt{1-\varkappa^2}$ gesetzt ist:

\[
\lambda = \frac{1-\varkappa'}{1+\varkappa'}, \qquad a=(1+\varkappa')^2,
\]

\[
\begin{aligned}
U-V &= [\xi-(1+\varkappa')\eta]^2,\\
U-\lambda^2V &= [\xi-(1-\varkappa')\eta]^2,\\
H &= (1+\varkappa')^2[\xi-(1+\varkappa')\eta][\xi-(1-\varkappa')\eta],\\
T &= (1+\varkappa')[(\xi-\eta)^2-\varkappa'^2\eta^2],\\
M &= \frac{1}{1+\varkappa'},
\end{aligned}
\]

also durch die Substitution

\weq{4}{z=\frac{(1+\varkappa')^2\zeta(1-\zeta)}{1-\varkappa^2\zeta}, \qquad
\lambda=\frac{1-\varkappa'}{1+\varkappa'}}

die Umformung:

\weq{5}{\frac{dz}{\sqrt{z(1-z)(1-\lambda^2z)}} =
\frac{(1+\varkappa')\,d\zeta}{\sqrt{\zeta(1-\zeta)(1-\varkappa^2\zeta)}}.}

Wendet man auf die linke Seite dieser Gleichung wieder die
\gesperrt{Gausssche} Transformation an, indem man in den Formeln (2), (3)
$\zeta,\varkappa$ durch $z,\lambda$ und $z$ durch eine Variable $\eta$
ersetzt, so ist, wie aus (2) und (4) folgt, $\lambda$ in (3) durch
$\varkappa$ zu ersetzen, und durch Kombination von (5) und (3) findet man:

\weq{6}{\frac{d\eta}{\sqrt{\eta(1-\eta)(1-\varkappa^2\eta)}} =
\frac{2\,d\zeta}{\sqrt{\zeta(1-\zeta)(1-\varkappa^2\zeta)}},}

und die Verbindung von (2) und (4) ergibt zwischen den Variablen $\eta$
und $\zeta$ den Zusammenhang:

\weq{7}{\eta = \frac{4\zeta(1-\zeta)(1-\varkappa^2\zeta)}{(1-\varkappa^2\zeta^2)^2}.}
